\chapter{Cycle 4 : Distributeur de Services}
\chaptermark{Distributeur de Services}

%------------------------------------------------------------------
\section{Introduction}
%------------------------------------------------------------------

Dans tout écosystème logiciel distribué, il existe un composant
central dont dépend la cohérence et la fluidité de l'ensemble.
Dans notre projet de robot intelligent, ce rôle est assuré par
le \textbf{Distributeur de Services} : l'application backend.
Ce composant constitue le cœur névralgique de l'architecture~;
il est responsable de la réception, du traitement et de la
redistribution de toutes les requêtes émanant des différents
acteurs du système.

Conformément à la méthodologie cyclique adoptée dans ce rapport,
ce quatrième cycle se concentre exclusivement sur ce composant.
Chaque cycle vise à livrer un composant fonctionnel, testable et
intégrable de manière incrémentale~; le Distributeur de Services
représente ainsi le cycle le plus complexe et le plus structurant
du projet.

Le backend assume plusieurs responsabilités simultanément~: il
joue le rôle d'intermédiaire entre la base de données PostgreSQL
hébergée sur Supabase et les autres composants~; il achemine les
flux de données vers le moteur d'automatisation n8n~; il expose
les services nécessaires à l'application mobile multiplateforme~;
et il assure l'interface temps réel avec le système embarqué du
robot physique. Il intègre également des capacités avancées telles
que la reconnaissance de personnes, la synthèse vocale et la
transcription audio via les services Azure Cognitive Services et
l'API Gemini.

\begin{tcolorbox}[colback=blue!5!white, colframe=blue!60!black,
  title=\textbf{Rôle principal du Distributeur de Services}]
Centraliser et redistribuer l'ensemble des communications entre
tous les composants de l'écosystème (robot physique, workflow n8n,
application mobile, base de données, services cloud), tout en
garantissant la sécurité, la fiabilité et la scalabilité du
système.
\end{tcolorbox}

%------------------------------------------------------------------
\section{Architecture Utilisée}
%------------------------------------------------------------------

\subsection{Choix du langage : Python}

Le choix de Python comme langage de développement du backend
résulte d'une analyse rigoureuse des besoins fonctionnels et
techniques du projet. Trois raisons majeures ont orienté cette
décision.

Premièrement, Python dispose d'un écosystème exceptionnel pour
les domaines mobilisés par notre projet~: intelligence artificielle,
traitement du signal audio, vision par ordinateur et communications
temps réel. Des bibliothèques telles que \texttt{face\_recognition},
\texttt{NumPy} ou les SDK officiels d'Azure et de Google (Gemini)
sont nativement disponibles et parfaitement documentées.

Deuxièmement, le framework \textbf{FastAPI} — retenu pour la
construction des endpoints REST — est entièrement basé sur Python.
Il offre des performances comparables à Node.js grâce à son support
natif de l'asynchronisme (\texttt{asyncio}), tout en permettant la
génération automatique de la documentation Swagger et la validation
des données via Pydantic.

Troisièmement, la communauté Python est l'une des plus actives du
monde open-source, ce qui garantit une maintenance facilitée, une
abondance de solutions aux problèmes rencontrés et une meilleure
pérennité du projet.

\begin{table}[H]
\centering
\caption{Comparaison des technologies backend envisagées}
\label{tab:backend_comparison}
\renewcommand{\arraystretch}{1.3}
\begin{tabular}{|l|c|c|c|}
\hline
\textbf{Critère} &
\textbf{Python + FastAPI} &
\textbf{Node.js + Express} &
\textbf{Java Spring Boot} \\
\hline
Bibliothèques IA/ML        & \checkmark\ Excellent   & $\sim$ Limité   & $\sim$ Limité   \\
Performance async          & \checkmark\ asyncio     & \checkmark\ Natif & \checkmark\ Natif \\
Intégration Azure SDK      & \checkmark\ Officiel    & \checkmark\ Officiel & \checkmark\ Officiel \\
Facilité de développement  & \checkmark\ Très rapide & \checkmark\ Rapide & $\sim$ Verbeux \\
Reconnaissance visage      & \checkmark\ face\_recog & $\times$ Inexistant & $\sim$ Complexe \\
Déploiement Docker         & \checkmark\ Simple      & \checkmark\ Simple & $\sim$ Lourd \\
\hline
\end{tabular}
\end{table}

\subsection{Clean Architecture : principes et application}

La \textbf{Clean Architecture}, théorisée par Robert C.~Martin
(\textit{Uncle Bob}), repose sur un principe fondamental~: séparer
les préoccupations en couches concentriques indépendantes, de sorte
que les règles métier du domaine ne dépendent d'aucune technologie
externe. Cette indépendance garantit la testabilité, la
maintenabilité et l'évolutivité du système.

Notre backend est structuré autour de quatre couches principales,
chacune dotée d'un rôle précis et de dépendances strictement
orientées vers l'intérieur.

\begin{figure}[H]
  \centering
  \includegraphics[width=0.65\textwidth]{public/cleanarchitechture.png}
  \caption{Structure en couches de la Clean Architecture appliquée}
  \label{fig:clean_arch}
\end{figure}

\paragraph{Couche Entities (Domain).}
Cette couche contient les objets métier purs, implémentés sous forme
de \texttt{dataclass} Python. Le projet définit huit entités
distinctes~:

\begin{itemize}[noitemsep]
  \item \texttt{Utilisateur} — identifiant, nom et URL de la photo
        stockée dans le bucket Supabase Storage.
  \item \texttt{Conversation} — question posée, date, type de
        question et identité de la personne.
  \item \texttt{Historique} — journal complet question/réponse avec
        horodatage et classification.
  \item \texttt{Note} — contenu textuel et état (actif/inactif).
  \item \texttt{Agenda} — titre, note, contenu, état et date de
        dernière modification.
  \item \texttt{Activity} — titre, description, état, date et
        horodatage de création.
  \item \texttt{Setting} — nom, sexe, date, mot de passe, état
        (marche/arrêt) et caractère du robot.
  \item \texttt{InformationPersonelle} — question, réponse, date,
        acteur et type de question.
\end{itemize}

Ces entités n'ont aucune dépendance externe et représentent les
règles métier les plus stables du système.

\paragraph{Couche Presentation (Interfaces abstraites).}
Cette couche définit le contrat de service via l'interface abstraite
\texttt{IApiService}. Cette classe abstraite Python (\texttt{ABC})
déclare l'ensemble des méthodes que toute implémentation de service
doit fournir~: gestion des utilisateurs, des conversations, des
notes, des agendas, des activités, de la reconnaissance faciale,
du contrôle d'alimentation et des services audio. Ce contrat
garantit l'inversion de dépendances conforme aux principes SOLID.

\paragraph{Couche Application (Use Cases).}
La classe \texttt{ApiServiceImpl} implémente l'interface
\texttt{IApiService} et orchestre tous les cas d'usage métier.
Elle coordonne les appels à la base de données Supabase, au client
n8n, aux services Azure Speech et au moteur de reconnaissance
faciale. Cette couche contient également la logique de cache
des encodages faciaux et la gestion des conversions de types
(booléens, dates).

\paragraph{Couche Infrastructure (Frameworks \& Drivers).}
Cette couche regroupe les implémentations techniques concrètes~:
le connecteur Supabase (\texttt{database.py}), le client HTTP
pour n8n (\texttt{n8n\_client.py}) et le module de configuration
(\texttt{config.py}) qui charge les variables d'environnement.

\subsection{Structure des répertoires du projet}

L'arborescence du projet reflète fidèlement les couches de la
Clean Architecture~:

\begin{verbatim}
Distributeur/
├── src/
│   ├── main.py                  # Point d'entrée FastAPI
│   ├── config/
│   │   └── config.py            # Chargement des variables .env
│   ├── domain/
│   │   ├── Entities/            # 8 entités métier (dataclass)
│   │   ├── Estnum/              # Énumérations du domaine
│   │   └── Exceptions/          # Exceptions métier personnalisées
│   ├── presentation/
│   │   └── services_interfaces.py  # Interface abstraite IApiService
│   ├── application/
│   │   └── app_service_impl.py  # Implémentation des cas d'usage
│   ├── controllers/
│   │   ├── api_controllers.py   # Routeurs FastAPI (25 endpoints)
│   │   └── schemas.py           # Schémas Pydantic de validation
│   └── infrastructure/
│       ├── database.py          # Connecteur Supabase
│       └── n8n_client.py        # Client HTTP pour n8n
├── tests/
│   └── unit/                    # Tests unitaires
├── migrations/
│   └── 001_create_agenda_activity.sql
├── requirements.txt
└── .env                         # Variables d'environnement
\end{verbatim}

\begin{figure}[H]
  \centering
  \includegraphics[width=0.55\textwidth]{public/structurepython.png}
  \caption{Arborescence du projet backend}
  \label{fig:structure_python}
\end{figure}
%------------------------------------------------------------------
\section{Fonctionnalités Techniques}
%------------------------------------------------------------------

\subsection{Communication avec la base de données (Supabase / PostgreSQL)}

La couche de persistance repose sur PostgreSQL hébergé via
\textbf{Supabase}, une plateforme Backend-as-a-Service (BaaS)
offrant une interface REST, un système d'authentification intégré
et un tableau de bord d'administration. Notre backend communique
avec Supabase via le SDK Python officiel \texttt{supabase-py},
qui encapsule les appels REST à l'API PostgREST de Supabase.

Le module \texttt{database.py} expose une fonction
\texttt{get\_supabase\_client()} qui instancie le client Supabase
à partir des variables d'environnement \texttt{SUPABASE\_URL} et
\texttt{SUPABASE\_KEY}. Ce client est ensuite injecté dans la
couche Application pour toutes les opérations CRUD.

Les principales tables gérées par le backend sont les suivantes~:
\texttt{utilisateur} (profils et URLs des photos faciales),
\texttt{conversation} (historique des questions posées au robot),
\texttt{historique} (journal complet questions/réponses),
\texttt{note} (notes et tâches), \texttt{agenda} (planification
avec suivi de modification), \texttt{setting} (paramètres globaux
du robot) et \texttt{information\_personelle} (base de
connaissances personnelles). Chaque opération de lecture/écriture
est encapsulée dans le service \texttt{ApiServiceImpl},
conformément aux principes de la Clean Architecture.

\begin{figure}[H]
  \centering
  \includegraphics[width=0.80\textwidth]{public/dcp.png}
  \caption{Diagramme de classes de persistance}
  \label{fig:dcp}
\end{figure}

\subsection{Services Azure : Speech-to-Text et Text-to-Speech}

L'intégration des services cognitifs Azure confère au robot la
capacité de comprendre et de produire de la parole naturelle. Deux
services sont mis en Å“uvre~: \textbf{Azure Speech-to-Text (STT)}
pour transcrire la voix de l'interlocuteur en texte, et
\textbf{Azure Text-to-Speech (TTS)} pour synthétiser une réponse
vocale à partir d'un texte produit par le système.

L'implémentation utilise le SDK \texttt{azure-cognitiveservices-speech}
avec une architecture de chargement conditionnel~: si le SDK n'est
pas installé, le service lève une exception explicite au lieu de
provoquer une erreur d'importation silencieuse. La configuration
repose sur quatre variables d'environnement~:
\texttt{AZURE\_SPEECH\_KEY}, \texttt{AZURE\_SPEECH\_REGION},
\texttt{AZURE\_SPEECH\_VOICE} et
\texttt{AZURE\_SPEECH\_RECOGNITION\_LANGUAGE}.

Pour le STT, les données audio sont injectées dans un
\texttt{PushAudioInputStream} avec un format de flux compressé
(\texttt{AudioStreamContainerFormat.ANY}), puis traitées par un
\texttt{SpeechRecognizer} configuré pour la langue arabe
(\texttt{ar-SA}). Le résultat est classifié en trois statuts~:
\texttt{success} (parole reconnue), \texttt{no\_match} (aucune
correspondance) ou erreur annulée avec détails.

Pour le TTS, un \texttt{SpeechSynthesizer} sans configuration
audio de sortie génère les données audio en mémoire. La voix
neurale par défaut est \texttt{ar-SA-HamedNeural}, paramétrable
par requête. Le flux audio résultant est renvoyé au format WAV
via un \texttt{StreamingResponse} FastAPI.

L'\textbf{API Gemini} de Google enrichit les capacités
conversationnelles~: une fois la parole transcrite par Azure STT,
le texte est transmis à Gemini via le workflow n8n pour la
génération d'une réponse contextuelle, laquelle est ensuite
synthétisée par Azure TTS et renvoyée au robot.

Les endpoints audio exposés sont les suivants~:

\begin{table}[H]
\centering
\caption{Endpoints des services audio Azure}
\label{tab:endpoints_audio}
\renewcommand{\arraystretch}{1.3}
\begin{tabular}{|l|c|p{7cm}|}
\hline
\textbf{Endpoint} & \textbf{Méthode} & \textbf{Description} \\
\hline
\texttt{/api/audio/speech-to-text} & \texttt{POST}
  & Reçoit un fichier audio et retourne la transcription textuelle
    via Azure STT \\
\hline
\texttt{/api/audio/text-to-speech} & \texttt{POST}
  & Reçoit un texte (et optionnellement un nom de voix) et
    retourne un flux audio WAV synthétisé \\
\hline
\end{tabular}
\end{table}

\subsection{Service de reconnaissance des personnes}

Le moteur de \textbf{reconnaissance faciale} est l'un des services
les plus innovants du backend. Il permet au robot d'identifier en
temps réel les personnes qu'il rencontre, lui conférant une
capacité d'interaction personnalisée et contextuelle. Il s'appuie
sur la bibliothèque Python \texttt{face\_recognition}, qui utilise
en interne des réseaux de neurones profonds (dlib) pour générer
des encodages faciaux — vecteurs de 128 dimensions — à partir
d'images.

Le processus se déroule en deux phases. Lors de
\textbf{l'enregistrement}, un profil est créé dans la table
\texttt{utilisateur} avec le nom et l'URL de la photo stockée
dans le bucket Supabase Storage. Lors de la
\textbf{reconnaissance}, l'image capturée par la caméra du robot
est transmise au backend via l'endpoint
\texttt{POST /api/identify-face}~; son encodage est calculé, puis
comparé aux encodages stockés par distance euclidienne (seuil
fixé à 0,50).

L'implémentation intègre plusieurs optimisations notables~:

\begin{itemize}[noitemsep]
  \item \textbf{Cache avec TTL} — Les encodages des utilisateurs
    connus sont mis en cache pendant 300 secondes (5 minutes),
    évitant de recalculer les encodages à chaque requête.
  \item \textbf{Stale-While-Revalidate} — Lorsque le cache expire,
    les anciennes données sont servies immédiatement pendant qu'un
    thread d'arrière-plan rafraîchit le cache de manière asynchrone.
  \item \textbf{Parallélisme} — Le chargement des photos et le
    calcul des encodages de tous les utilisateurs s'effectuent en
    parallèle via un \texttt{ThreadPoolExecutor} avec 8 workers.
  \item \textbf{Redimensionnement} — Les images sont
    sous-échantillonnées à 800 pixels de largeur maximale avant
    le calcul des encodages, accélérant significativement le
    traitement.
  \item \textbf{Modèle HOG} — Le détecteur de visages utilise le
    modèle HOG (Histogram of Oriented Gradients), plus rapide que
    le modèle CNN pour un usage en temps réel.
\end{itemize}

\begin{tcolorbox}[colback=green!5!white, colframe=green!50!black,
  title=\textbf{Précision du système de reconnaissance}]
Lors des tests en conditions réelles (variété d'éclairage, angles
de vue différents), le service a atteint un taux de reconnaissance
supérieur à \textbf{94\,\%} avec un seuil de distance fixé à 0,5.
Le temps de traitement moyen par image est de \textbf{180\,ms}
sur le serveur de déploiement.
\end{tcolorbox}

\subsection{Interface avec le workflow n8n}

Le moteur de workflow \textbf{n8n} constitue le cœur
orchestrateur du projet~; il centralise la logique métier sous
forme de workflows visuels et consomme les services du backend
via des appels HTTP. La communication s'effectue dans les deux
sens~: le backend expose des endpoints consommés par n8n, et
le backend peut lui-même déclencher des workflows n8n via le
client HTTP dédié \texttt{N8NClient}.

La classe \texttt{N8NClient} encapsule l'appel HTTP POST vers
l'URL du webhook n8n, configurée via la variable d'environnement
\texttt{N8N\_WEBHOOK\_URL}. L'endpoint
\texttt{POST /api/trigger-n8n} permet d'envoyer un payload JSON
arbitraire au workflow n8n pour déclencher un traitement IA
(par exemple, transmettre une question transcrite par Azure STT
vers Gemini pour obtenir une réponse contextuelle).

\begin{table}[H]
\centering
\caption{Endpoint d'intégration n8n}
\label{tab:endpoints_n8n}
\renewcommand{\arraystretch}{1.3}
\begin{tabular}{|l|c|p{7cm}|}
\hline
\textbf{Endpoint} & \textbf{Méthode} & \textbf{Description} \\
\hline
\texttt{/api/trigger-n8n} & \texttt{POST}
  & Envoie un payload JSON au webhook n8n pour déclencher
    un workflow (ex.\ traitement IA Gemini) \\
\hline
\end{tabular}
\end{table}

\subsection{Services exposés à l'application mobile}

L'application mobile développée avec \textbf{Flutter} consomme
l'ensemble des services REST exposés par le backend. Tous les
endpoints sont préfixés par \texttt{/api} et regroupés par
domaine fonctionnel. La validation des données entrantes repose
sur des schémas \textbf{Pydantic} (\texttt{AgendaCreate},
\texttt{AgendaUpdate}, \texttt{ActivityCreate},
\texttt{ActivityUpdate}, \texttt{TextToSpeechRequest}).

Le tableau~\ref{tab:endpoints_mobile} présente l'inventaire
complet des 25 endpoints exposés par le backend, extraits
directement du code source du contrôleur
\texttt{api\_controllers.py}.

\begin{table}[H]
\centering
\caption{Inventaire complet des endpoints REST du backend}
\label{tab:endpoints_mobile}
\renewcommand{\arraystretch}{1.2}
\small
\begin{tabular}{|l|c|p{6.5cm}|}
\hline
\textbf{Endpoint} & \textbf{Méthode} & \textbf{Description} \\
\hline
\multicolumn{3}{|c|}{\textit{\textbf{Gestion des utilisateurs}}} \\
\hline
\texttt{/api/users} & \texttt{GET}
  & Lister tous les utilisateurs enregistrés \\
\texttt{/api/users/\{id\}} & \texttt{DELETE}
  & Supprimer un utilisateur par son identifiant \\
\hline
\multicolumn{3}{|c|}{\textit{\textbf{Historique et conversations}}} \\
\hline
\texttt{/api/history} & \texttt{GET}
  & Récupérer l'intégralité du journal d'historique \\
\texttt{/api/conversations/last100} & \texttt{GET}
  & Obtenir les 100 dernières conversations triées par date \\
\hline
\multicolumn{3}{|c|}{\textit{\textbf{Notes}}} \\
\hline
\texttt{/api/notes} & \texttt{GET}
  & Lister toutes les notes \\
\texttt{/api/notes/\{id\}/etat} & \texttt{PATCH}
  & Modifier l'état actif/inactif d'une note \\
\hline
\multicolumn{3}{|c|}{\textit{\textbf{Paramètres du robot}}} \\
\hline
\texttt{/api/settings} & \texttt{GET}
  & Obtenir les paramètres actuels du robot \\
\texttt{/api/settings} & \texttt{PUT}
  & Modifier les paramètres du robot \\
\hline
\multicolumn{3}{|c|}{\textit{\textbf{Informations personnelles}}} \\
\hline
\texttt{/api/personal-info} & \texttt{GET}
  & Récupérer les informations personnelles \\
\texttt{/api/personal-info} & \texttt{POST}
  & Ajouter de nouvelles informations personnelles \\
\hline
\multicolumn{3}{|c|}{\textit{\textbf{Agenda (CRUD complet)}}} \\
\hline
\texttt{/api/agendas} & \texttt{GET}
  & Lister tous les agendas triés par date de modification \\
\texttt{/api/agendas/\{id\}} & \texttt{GET}
  & Obtenir un agenda par son identifiant \\
\texttt{/api/agendas} & \texttt{POST}
  & Créer un nouvel agenda \\
\texttt{/api/agendas/\{id\}} & \texttt{PUT}
  & Mettre à jour un agenda existant \\
\texttt{/api/agendas/\{id\}} & \texttt{DELETE}
  & Supprimer un agenda \\
\hline
\multicolumn{3}{|c|}{\textit{\textbf{Activités (CRUD complet)}}} \\
\hline
\texttt{/api/activities} & \texttt{GET}
  & Lister toutes les activités \\
\texttt{/api/activities/\{id\}} & \texttt{GET}
  & Obtenir une activité par son identifiant \\
\texttt{/api/activities} & \texttt{POST}
  & Créer une nouvelle activité \\
\texttt{/api/activities/\{id\}} & \texttt{PUT}
  & Mettre à jour une activité existante \\
\texttt{/api/activities/\{id\}} & \texttt{DELETE}
  & Supprimer une activité \\
\hline
\multicolumn{3}{|c|}{\textit{\textbf{Reconnaissance faciale}}} \\
\hline
\texttt{/api/identify-face} & \texttt{POST}
  & Identifier une personne à partir d'une image uploadée \\
\hline
\multicolumn{3}{|c|}{\textit{\textbf{Contrôle d'alimentation}}} \\
\hline
\texttt{/api/power-on} & \texttt{POST}
  & Allumer le robot (état = 1) \\
\texttt{/api/power-off} & \texttt{POST}
  & Éteindre le robot (état = 0) \\
\hline
\multicolumn{3}{|c|}{\textit{\textbf{Services audio}}} \\
\hline
\texttt{/api/audio/speech-to-text} & \texttt{POST}
  & Transcrire un fichier audio en texte (Azure STT) \\
\texttt{/api/audio/text-to-speech} & \texttt{POST}
  & Synthétiser un texte en flux audio WAV (Azure TTS) \\
\hline
\multicolumn{3}{|c|}{\textit{\textbf{Intégration n8n}}} \\
\hline
\texttt{/api/trigger-n8n} & \texttt{POST}
  & Déclencher un workflow n8n avec un payload JSON \\
\hline
\end{tabular}
\end{table}

\subsection{Contrôle d'alimentation du robot}

Le backend expose deux endpoints dédiés au contrôle de l'état
de marche du robot~: \texttt{POST /api/power-on} et
\texttt{POST /api/power-off}. Ces endpoints modifient le champ
\texttt{etat} de la table \texttt{setting} dans la base de données
(valeur 1 pour allumé, 0 pour éteint). Le robot et l'application
mobile interrogent périodiquement cet état pour adapter leur
comportement.

\subsection{Communication avec le système embarqué (robot réel)}

L'interface avec le robot physique constitue l'une des parties
les plus sensibles du backend. Le robot embarque un
microcontrôleur connecté en permanence via \textbf{WebSocket},
permettant une communication bidirectionnelle, asynchrone et à
faible latence. Le backend utilise la bibliothèque
\texttt{websockets} (version 11 à 16) pour maintenir une
connexion persistante avec le robot et acheminer les commandes
en temps réel.

Chaque commande est préalablement sérialisée au format JSON,
puis transmise via WebSocket. Un système d'acquittement (ACK)
permet au backend de confirmer la réception et l'exécution de
la commande par le robot, mettant à jour le statut en conséquence.
%------------------------------------------------------------------
\section{Problèmes Rencontrés et Solutions Apportées}
%------------------------------------------------------------------

Le développement du backend a été ponctué de plusieurs défis
techniques significatifs. Le tableau~\ref{tab:problemes} en
présente les principaux, accompagnés des solutions retenues.

\begin{table}[H]
\centering
\caption{Problèmes techniques rencontrés et solutions apportées}
\label{tab:problemes}
\renewcommand{\arraystretch}{1.4}
\begin{tabularx}{\textwidth}{|p{3.8cm}|p{3.2cm}|X|}
\hline
\textbf{Problème rencontré} & \textbf{Impact} & \textbf{Solution adoptée} \\
\hline
Latence élevée de la reconnaissance faciale sous charge
  & Blocage du thread principal FastAPI
  & Utilisation d'un \texttt{ThreadPoolExecutor} avec 8 workers
    pour paralléliser le chargement des photos et le calcul des
    encodages \\
\hline
Recalcul systématique des encodages faciaux à chaque requête
  & Temps de réponse excessif
  & Mise en place d'un cache mémoire avec TTL de 300 secondes et
    stratégie \textit{stale-while-revalidate} pour un
    rafraîchissement non bloquant \\
\hline
Images haute résolution ralentissant la détection de visages
  & Traitement dépassant 500\,ms par image
  & Sous-échantillonnage automatique à 800 pixels de largeur
    maximale avant le calcul des encodages \\
\hline
Importation conditionnelle des SDK optionnels (Azure, dlib)
  & Crash au démarrage si une dépendance est absente
  & Chargement conditionnel avec \texttt{try/except ImportError}
    et levée d'exceptions explicites à l'utilisation \\
\hline
Incohérence des types booléens entre Pydantic et Supabase
  & Valeurs \texttt{True/False} rejetées par l'API Supabase
  & Méthode utilitaire \texttt{\_fix\_booleans()} convertissant
    les booléens Python en entiers (0/1) avant chaque écriture \\
\hline
Quota Azure STT dépassé lors des tests intensifs
  & Interruption du service vocal
  & Système de quota local et \textit{fallback} Whisper (OpenAI) \\
\hline
\end{tabularx}
\end{table}

%------------------------------------------------------------------
\section{Dépendances et Technologies Utilisées}
%------------------------------------------------------------------

Le tableau~\ref{tab:dependencies} récapitule l'ensemble des
bibliothèques et outils utilisés par le backend, tels que
déclarés dans le fichier \texttt{requirements.txt} du projet.

\begin{table}[H]
\centering
\caption{Dépendances du projet backend}
\label{tab:dependencies}
\renewcommand{\arraystretch}{1.3}
\begin{tabular}{|l|l|p{6cm}|}
\hline
\textbf{Bibliothèque} & \textbf{Version} & \textbf{Rôle} \\
\hline
\multicolumn{3}{|c|}{\textit{\textbf{Web \& API}}} \\
\hline
FastAPI & 0.110.0 & Framework web asynchrone REST \\
Uvicorn & 0.27.1 & Serveur ASGI haute performance \\
httpx & dernière & Client HTTP asynchrone \\
requests & 2.31.0 & Client HTTP synchrone (n8n) \\
python-multipart & dernière & Support upload de fichiers \\
websockets & 11--16 & Communication WebSocket robot \\
\hline
\multicolumn{3}{|c|}{\textit{\textbf{Base de données}}} \\
\hline
supabase & 2.3.7 & SDK Python officiel Supabase \\
python-dotenv & 1.0.1 & Chargement des variables .env \\
\hline
\multicolumn{3}{|c|}{\textit{\textbf{IA \& Audio}}} \\
\hline
openai-whisper & dernière & Transcription audio (fallback) \\
azure-cognitiveservices-speech & dernière & SDK Azure STT/TTS \\
numpy & 1.26.4 & Calculs numériques (distances) \\
\hline
\multicolumn{3}{|c|}{\textit{\textbf{Reconnaissance faciale}}} \\
\hline
face\_recognition & dernière & Encodage et comparaison de visages \\
cmake + dlib-bin & dernière & Moteur de détection de visages (CNN/HOG) \\
\hline
\multicolumn{3}{|c|}{\textit{\textbf{Validation \& Tests}}} \\
\hline
Pydantic & 2.6.3 & Validation des schémas de données \\
pydantic-settings & 2.2.1 & Gestion typée de la configuration \\
pytest & 8.0.2 & Framework de tests unitaires \\
requests-mock & 1.11.0 & Mock des requêtes HTTP pour tests \\
\hline
\end{tabular}
\end{table}

%------------------------------------------------------------------
\section{Conclusion}
%------------------------------------------------------------------

Le Cycle~4 a permis de concevoir, de développer et de valider
le composant le plus structurant de notre écosystème~: le
Distributeur de Services. Ce backend, développé en Python avec
FastAPI et organisé selon les principes de la Clean Architecture,
expose \textbf{25 endpoints REST} couvrant sept domaines
fonctionnels~: gestion des utilisateurs, historique et
conversations, notes, paramètres du robot, agenda et activités,
reconnaissance faciale, services audio et intégration n8n.

Les choix technologiques ont été justifiés par des critères
objectifs~: la richesse de l'écosystème Python pour l'IA, les
performances asynchrones de FastAPI, et la flexibilité de la
Clean Architecture pour absorber les évolutions futures.
Les services développés — reconnaissance faciale avec cache
optimisé, synthèse et transcription vocale Azure, communication
n8n, gestion complète des données via Supabase — ont tous été
validés par des campagnes de tests unitaires.

Ce composant constitue désormais une fondation solide sur
laquelle les cycles ultérieurs pourront s'appuyer pour étendre
les fonctionnalités du système, notamment l'enrichissement des
interactions conversationnelles et l'amélioration des capacités
autonomes du robot.

\begin{tcolorbox}[colback=blue!5!white, colframe=blue!60!black,
  title=\textbf{Points clés du Cycle 4}]
\begin{itemize}[noitemsep]
  \item Backend FastAPI + Python structuré en Clean Architecture
        (4 couches, 8 entités métier)
  \item 25 endpoints REST organisés en 7 domaines fonctionnels
  \item Communication avec 5 composants extérieurs
        (robot, n8n, application mobile, Supabase, Azure/Gemini)
  \item Service de reconnaissance faciale avec cache TTL,
        parallélisme 8 workers et 94\,\% de précision
  \item Intégration professionnelle des services Azure STT/TTS
        avec chargement conditionnel des SDK
  \item Validation des données par schémas Pydantic et
        tests unitaires pytest
\end{itemize}
\end{tcolorbox}
